% SGA 3, Expose I, Loop 1 English reconstruction.
% Scope: Section 2.3, from the heading through Proposition 2.3.8 and
% editor's notes (23)--(26); stops immediately before Section 3.
% Authority: Polo--Gille Exp1-13oct24.pdf, component local pages 15--17
% (combined-reader pages 25--27). OCR is drafting/locator material only.

\subsection*{\SGAInlineRefTarget{sga3:i:subsection:2-3}2.3. 带算子群的结构}

\medskip
\noindent\SGAInlineRefTarget{sga3:i:definition:2-3-1}\textbf{定义 2.3.1.}
设 \(\mathbf E\in\operatorname{Ob}(\widehat{\mathscr C})\) 且
\(\mathbf G\in
\operatorname{Ob}(\widehat{\mathscr C}\text{-}\mathbf{Gr})\).
在 \(\mathbf E\) 上赋予以 \(\widehat{\mathscr C}\)-群 \(\mathbf G\)
为算子群的对象结构 (或称 \(\mathbf G\)-对象结构), 是指对每个
\(S\in\operatorname{Ob}(\mathscr C)\), 在 \(\mathbf E(S)\) 上赋予以
\(\mathbf G(S)\) 为算子群的集合结构, 使得对于 \(\mathscr C\) 中的每个箭头
\(S'\to S''\), 集合映射
\(\mathbf E(S'')\to\mathbf E(S')\) 与算子群同态
\(\mathbf G(S'')\to\mathbf G(S')\) 相容.

照常, 这等价于给定一个态射
\[
  \mu:\mathbf G\times\mathbf E\longrightarrow\mathbf E
\]
使其对每个 \(S\) 都在 \(\mathbf E(S)\) 上赋予一个
\(\mathbf G(S)\) 作用的集合结构. 然而,
\[
  \operatorname{Hom}(\mathbf G\times\mathbf E,\mathbf E)
    \simeq
  \operatorname{Hom}
    (\mathbf G,\underline{\operatorname{End}}(\mathbf E)),
\]
所以 \(\mu\) 定义了一个态射
\(\mathbf G\to\underline{\operatorname{End}}(\mathbf E)\). 显然, 这个态射将
\(\mathbf G\) 映入 \(\underline{\operatorname{Aut}}(\mathbf E)\),
并且是一个 \(\widehat{\mathscr C}\)-群态射. 因此, 在 \(\mathbf E\) 上赋予
以 \(\widehat{\mathscr C}\)-群 \(\mathbf G\) 为算子群的对象结构,
等价于给定一个 \(\widehat{\mathscr C}\)-群态射
\[
  \rho:\mathbf G\longrightarrow
    \underline{\operatorname{Aut}}(\mathbf E).
\]
特别地, 每个元素 \(g\in\mathbf G(S)\) 都定义了函子 \(\mathbf E_S\)
的一个自同构 \(\rho(g)\), 即 \(\mathbf E\times\mathbf h_S\) 的一个
与投影 \(\mathbf E\times\mathbf h_S\to\mathbf h_S\) 可交换的自同构;
因而特别地, 对每个 \(S'\to S\), 都定义了集合 \(\mathbf E(S')\)
的一个自同构.

\medskip
\noindent\SGAInlineRefTarget{sga3:i:definition:2-3-2}\textbf{定义 2.3.2.}
以 \(\mathbf E^{\mathbf G}\) 表示 \(\mathbf E\) 的如下子对象:
\[
  \mathbf E^{\mathbf G}(S)
    =\bigl\{x\in\mathbf E(S)\mid
      x_{S'}\text{ 在 }\mathbf G(S')
      \text{ 的作用下不变, 对每个 }S'\to S\bigr\},
\]
其中 \(x_{S'}\) 表示 \(x\) 在
\(\mathbf E(S)\to\mathbf E(S')\) 下的像.

于是, \(\mathbf E^{\mathbf G}\), 即 ``\(\mathbf G\) 的不变量子对象'',
是 \(\mathbf E\) 的一个最大子对象, 其上 \(\mathbf G\) 平凡作用.

\medskip
\SGARefTarget{sga3:i:definition:2:3:3}
\noindent\textbf{定义 2.3.3.}
设 \(\mathbf F\) 是 \(\mathbf E\) 的一个子对象. 以
\(\underline{\operatorname{Norm}}_{\mathbf G}\mathbf F\) 和
\(\underline{\operatorname{Centr}}_{\mathbf G}\mathbf F\) 表示由下式定义的
\(\mathbf G\) 的子 \(\widehat{\mathscr C}\)-群:
\begin{align*}
  (\underline{\operatorname{Norm}}_{\mathbf G}\mathbf F)(S)
    &={}
      \{g\in\mathbf G(S)\mid \rho(g)\mathbf F_S=\mathbf F_S\}\\
    &={}
      \{g\in\mathbf G(S)\mid
        \rho(g)\mathbf F(S')=\mathbf F(S')
        \text{ 对每个 }S'\to S\},\\
  (\underline{\operatorname{Centr}}_{\mathbf G}\mathbf F)(S)
    &={}
      \{g\in\mathbf G(S)\mid
        \rho(g)|_{\mathbf F_S}=\operatorname{id}\}\\
    &={}
      \{g\in\mathbf G(S)\mid
        \rho(g)|_{\mathbf F(S')}=\operatorname{id}
        \text{ 对每个 }S'\to S\}.
\end{align*}
这里, \(\rho(g)\) 后的竖线表示限制.
\footnote{编者注: 原文已作修正, 将包含关系
\(\rho(g)\mathbf F_S\subset\mathbf F_S\) 改为等式, 以保证
\(\underline{\operatorname{Norm}}_{\mathbf G}(\mathbf F)\) 确实是一个群.
关于 ``输运子 (transporter)'' 与 ``严格输运子 (strict transporter)''
重合的条件, 见 VI\textsubscript{B}, \SGARefLink{prop:sga3-VIB-6.4}{6.4}.}

\medskip
\noindent\SGAInlineRefTarget{sga3:i:scholium:2-3-3-1}\textbf{附论 2.3.3.1.}
\footnote{编者注: \SGARefLink{sga3:i:scholium:2-3-3-1}{附论 2.3.3.1} 和
\SGARefLink{sga3:i:remark:2-3-3-2}{注 2.3.3.2} 为新增内容.}
特别地, 设 \(x\in\Gamma(\mathbf E)\), 即 (参见 \SGARefLink{sga3:i:subsection:1-2}{1.2})
一族元素 \(x_S\in\mathbf E(S)\),
\(S\in\operatorname{Ob}(\mathscr C)\), 使得对于每个箭头
\(f:S'\to S\), 都有
\(\mathbf E(f)(x_S)=x_{S'}\). 若 \(\mathscr C\) 有终对象
\(S_0\), 则 \(\Gamma(\mathbf E)=\mathbf E(S_0)\). 元素 \(x\)
于是定义了 \(\mathbf E\) 的一个子函子, 记作 \(\mathbf x\), 并且
\[
  \underline{\operatorname{Norm}}_{\mathbf G}\mathbf x
    =
  \underline{\operatorname{Centr}}_{\mathbf G}\mathbf x.
\]
将这个函子记作
\(\underline{\operatorname{Stab}}_{\mathbf G}(x)\), 并称其为 \(x\) 的
\emph{稳定子 (stabilizer)}; 于是, 对每个
\(S\in\operatorname{Ob}(\mathscr C)\),
\[
  \underline{\operatorname{Stab}}_{\mathbf G}(x)(S)
    =\{g\in\mathbf G(S)\mid\rho(g)x_S=x_S\}.
\]

假设 \(\mathscr C\) 中存在纤维积. 若
\(\mathbf G=\mathbf h_G\) (分别地,
\(\mathbf E=\mathbf h_E\)), 其中 \(G\) 是一个 \(\mathscr C\)-群
(分别地, \(E\in\operatorname{Ob}(\mathscr C)\)), 且
\(\mathscr C\) 有终对象 \(S_0\), 从而 \(x\) 是一个态射
\(S_0\to E\), 则
\(\underline{\operatorname{Stab}}_{\mathbf G}(x)\) 由纤维积
\(G\times_E S_0\) 表出, 其中态射 \(G\to E\) 是
\[
  \operatorname{id}_G\times x:
  G=G\times S_0\longrightarrow G\times E
\]
与
\[
  \mu:G\times E\longrightarrow E.
\]
的复合.

\medskip
\noindent\SGAInlineRefTarget{sga3:i:remark:2-3-3-2}\textbf{注 2.3.3.2.}\footnotemark[\value{footnote}]
\(\mathbf E^{\mathbf G}\),
\(\underline{\operatorname{Norm}}_{\mathbf G}\mathbf F\) 和
\(\underline{\operatorname{Centr}}_{\mathbf G}\mathbf F\) 的构造与基变换可交换;
即对每个
\(S\in\operatorname{Ob}(\mathscr C)\),
\[
  (\mathbf E^{\mathbf G})_S=(\mathbf E_S)^{\mathbf G_S},
  \qquad
  (\underline{\operatorname{Norm}}_{\mathbf G}\mathbf F)_S
    \simeq
  \underline{\operatorname{Norm}}_{\mathbf G_S}\mathbf F_S,
  \qquad
  (\underline{\operatorname{Centr}}_{\mathbf G}\mathbf F)_S
    \simeq
  \underline{\operatorname{Centr}}_{\mathbf G_S}\mathbf F_S.
\]

\medskip
\noindent\SGAInlineRefTarget{sga3:i:definition:2-3-4}\textbf{定义 2.3.4.}
若 \(G\) 是一个 \(\mathscr C\)-群, \(\mathbf E\) 是
\(\widehat{\mathscr C}\) 的对象 (分别地, \(E\) 是
\(\mathscr C\) 的对象), 则 \(\mathbf E\) 上 (分别地, \(E\) 上)
的 \(G\)-对象结构, 是指 \(\mathbf E\) 上 (分别地, \(\mathbf h_E\) 上)
的 \(\mathbf h_G\)-对象结构.

根据这个定义, 以上定义的所有概念和记号, 只要仅涉及可表函子,
都可移用到 \(\mathscr C\) 中. 例如, 若
\(\underline{\operatorname{Norm}}_{\mathbf h_G}(\mathbf h_F)\) 可表,
则有唯一的 \(G\) 的子对象表出它; 该子对象于是成为 \(G\) 的一个
子 \(\mathscr C\)-群, 记作 \(\operatorname{Norm}_G(F)\), 其余类推.

\medskip
\noindent\SGAInlineRefTarget{sga3:i:definition:2-3-5}\textbf{定义 2.3.5.}
\begin{enumerate}
  \item[(a)]
  若 \(\widehat{\mathscr C}\)-群 \(\mathbf H\) 配备了一个
  \(\mathbf G\)-对象结构, 并且对每个 \(g\in\mathbf G(S)\),
  由 \(g\) 定义的 \(\mathbf H(S)\) 的自同构都是群自同构,
  则称 \(\widehat{\mathscr C}\)-群 \(\mathbf G\) 作用于
  群 \(\mathbf H\).

  等价地, 对每个 \(g\in\mathbf G(S)\), \(\mathbf H_S\) 的自同构
  \(\rho(g)\) 是一个 \(\widehat{\mathscr C}_{/S}\)-群自同构;
  或者等价地, \(\widehat{\mathscr C}\)-群态射
  \(\mathbf G\to\underline{\operatorname{Aut}}(\mathbf H)\) 将
  \(\mathbf G\) 映入
  \[
    \underline{\operatorname{Aut}}_{
      \widehat{\mathscr C}\text{-}\mathbf{Gr}}(\mathbf H).
  \]

  \item[(b)]
  在上述情形中, \(\mathbf H\times\mathbf G\) 上存在唯一的
  \(\widehat{\mathscr C}\)-群结构, 使得对每个 \(S\),
  \((\mathbf H\times\mathbf G)(S)\) 都是群 \(\mathbf H(S)\) 与
  \(\mathbf G(S)\) 关于给定的 \(\mathbf G(S)\) 在
  \(\mathbf H(S)\) 上的作用的半直积. 将这个
  \(\widehat{\mathscr C}\)-群记作 \(\mathbf H\mathbin{\cdot}\mathbf G\),
  称为\emph{\(\mathbf H\) 与 \(\mathbf G\) 的半直积 (semidirect product)}.
  于是, 按定义,
  \[
    (\mathbf H\mathbin{\cdot}\mathbf G)(S)
      =\mathbf H(S)\mathbin{\cdot}\mathbf G(S).
  \]
\end{enumerate}

设 \(\mathbf G\) 是一个 \(\widehat{\mathscr C}\)-群.
对于 \(\mathscr C\) 中的每个箭头 \(S'\to S\) 以及每个
\(g\in\mathbf G(S)\), 以 \(\operatorname{Int}(g)\) 表示由下式定义的
\(\mathbf G(S')\) 的自同构:
\[
  \operatorname{Int}(g)h=ghg^{-1}.
\]
这个定义延拓为一个 \(\widehat{\mathscr C}\)-群态射
\[
  \operatorname{Int}:\mathbf G\longrightarrow
  \underline{\operatorname{Aut}}_{
    \widehat{\mathscr C}\text{-}\mathbf{Gr}}(\mathbf G)
  \subset
  \underline{\operatorname{Aut}}(\mathbf G).
\]

因此可应用\SGARefLink{sga3:i:definition:2:3:3}{定义 2.3.3};
对于 \(\mathbf G\) 的每个子对象 \(\mathbf E\), 都有 \(\mathbf G\)
的子 \(\widehat{\mathscr C}\)-群
\[
  \underline{\operatorname{Norm}}_{\mathbf G}(\mathbf E)
  \qquad\text{及}\qquad
  \underline{\operatorname{Centr}}_{\mathbf G}(\mathbf E).
\]

\medskip
\SGARefTarget{sga3:i:definition:2:3:6}
\noindent\textbf{定义 2.3.6.}
\(\mathbf G\) 的\emph{中心 (center)} 记作
\(\underline{\operatorname{Centr}}(\mathbf G)\), 是指 \(\mathbf G\)
的子 \(\widehat{\mathscr C}\)-群
\(\underline{\operatorname{Centr}}_{\mathbf G}(\mathbf G)\).
若 \(\underline{\operatorname{Centr}}(\mathbf G)=\mathbf G\),
或等价地, 对每个 \(S\), \(\mathbf G(S)\) 都是交换群,
则称 \(\mathbf G\) 是交换的.

若 \(\underline{\operatorname{Norm}}_{\mathbf G}(\mathbf H)=\mathbf G\),
或等价地, 对每个 \(S\), \(\mathbf H(S)\) 都是 \(\mathbf G(S)\)
的正规子群, 则称 \(\mathbf G\) 的子 \(\widehat{\mathscr C}\)-群
\(\mathbf H\) 在 \(\mathbf G\) 中正规.
\footnote{编者注: 此外, 若
\(\underline{\operatorname{Centr}}_{\mathbf G}(\mathbf H)=\mathbf G\),
或等价地, 对每个 \(S\), \(\mathbf H(S)\) 都包含于 \(\mathbf G(S)\)
的中心, 则称 \(\mathbf H\) 在 \(\mathbf G\) 中是\emph{中心的 (central)}.}

\medskip
\noindent\SGAInlineRefTarget{sga3:i:definition:2-3-6-1}\textbf{定义 2.3.6.1.}
\footnote{编者注: \SGARefLink{sga3:i:definition:2-3-6-1}{定义 2.3.6.1} 和
\SGARefLink{sga3:i:definition:2-3-6-2}{2.3.6.2} 为新增内容.}
设 \(f:\mathbf G\to\mathbf G'\) 是一个 \(\widehat{\mathscr C}\)-群态射.
\(f\) 的\emph{核 (kernel)} 记作 \(\operatorname{Ker}f\), 是指对每个
\(S\in\operatorname{Ob}(\mathscr C)\) 由下式定义的 \(\mathbf G\)
的子 \(\widehat{\mathscr C}\)-群:
\[
  (\operatorname{Ker}f)(S)
    =\{x\in\mathbf G(S)\mid f(S)(x)=1\}
    =\operatorname{Ker}f(S)
\]
它是一个正规子 \(\widehat{\mathscr C}\)-群.

注意, 若 \(\mathbf G=\mathbf h_G\),
\(\mathbf G'=\mathbf h_{G'}\), \(\mathscr C\) 有终对象
\(S_0\), 且 \(\mathscr C\) 中存在纤维积, 则
\(\operatorname{Ker}(f)\) 由 \(S_0\times_{G'}G\) 表出.

\medskip
\noindent\SGAInlineRefTarget{sga3:i:definition:2-3-6-2}\textbf{定义 2.3.6.2.}\footnotemark[\value{footnote}]
设 \(\mathbf E\in\widehat{\mathscr C}\), 而 \(\mathbf G\) 是一个作用于
\(\mathbf E\) 的 \(\widehat{\mathscr C}\)-群.
若态射 \(\mathbf G\to\underline{\operatorname{Aut}}(\mathbf E)\)
的核平凡, 则称 \(\mathbf G\) 在 \(\mathbf E\) 上的作用是
\emph{忠实的 (faithful)}; 即对每个
\(S\in\operatorname{Ob}(\mathscr C)\) 及 \(g\in\mathbf G(S)\), 条件
\[
  g_{S'}\cdot x=x
  \qquad
  \text{对每个 }S'\to S\text{ 及每个 }x\in\mathbf E(S')
\]
蕴含 \(g=1\).

初等群论中的许多定义和命题都很容易移用于这里.
我们只记下以下将要用到的命题.

\medskip
\noindent\SGAInlineRefTarget{sga3:i:proposition:2-3-7}\textbf{命题 2.3.7.}
\textit{设 \(f:\mathbf W\to\mathbf G\) 是一个
\(\widehat{\mathscr C}\)-群态射. 置
\(\mathbf H(S)=\operatorname{Ker}f(S)\). 设
\(u:\mathbf G\to\mathbf W\) 是一个作为 \(f\) 的截面的
\(\widehat{\mathscr C}\)-群态射 (所以必为单态射).
则 \(\mathbf W\) 可等同于 \(\mathbf H\) 与 \(\mathbf G\)
关于如下 \(\mathbf G\) 在 \(\mathbf H\) 上的作用的半直积:
\[
  (g,h)\longmapsto\operatorname{Int}(u(g))h
\]
其中 \(g\in\mathbf G(S)\), \(h\in\mathbf H(S)\),
\(S\in\operatorname{Ob}(\mathscr C)\).}

照常, 所有这些定义和命题都可移用于 \(\mathscr C\).
特别地, 若 \(G\) 作用于 \(H\), 则只要笛卡尔积 \(H\times G\) 存在,
两个 \(\mathscr C\)-群 \(H\) 与 \(G\) 的半直积便有定义;
而\SGARefLink{sga3:i:proposition:2-3-7}{命题 2.3.7} 的类似命题如下.

\medskip
\noindent\SGAInlineRefTarget{sga3:i:proposition:2-3-8}\textbf{命题 2.3.8.}
\textit{设
\[
  H\xrightarrow{i}W\xrightarrow{f}G
\]
是一个 \(\mathscr C\)-群态射序列, 使得对每个
\(S\in\operatorname{Ob}(\mathscr C)\), \((H(S),i(S))\) 都是
\(f(S):W(S)\to G(S)\) 的一个核. 设 \(u:G\to W\) 是一个
作为 \(f\) 的截面的 \(\mathscr C\)-群态射.
则 \(W\) 可等同于 \(H\) 与 \(G\) 关于如下 \(G\) 在 \(H\)
上的作用的半直积: 若 \(S\in\operatorname{Ob}(\mathscr C)\),
\(g\in G(S)\), \(h\in H(S)\), 则
\[
  \operatorname{Int}(u(g))i(h)=i({}^{g}h).
\]}

% Translation coverage: entire assigned component; no unresolved translation TODOs.
% Statistics: source lines=346; target lines=328; inline math=223; display math=19; align blocks=1; unresolved TODOs=0.
